% Owner: Hassan Ashraf 24L-3073
\section{Buy, Hold and Negotiate Recommendations}
\label{sec:recommendation}

The models in Sections~\ref{sec:classical} and~\ref{sec:deep} predict only the
direction of gold. A buyer, however, needs an answer to a different question:
\emph{is this supplier's price good, and should I buy now?} This section
describes how ProcurePulse turns market data and the model output into one of
three actions for a supplier's quote. The code is in
\texttt{src/recommendation/price\_range.py} and is used by the Quote Check page
of the application (Section~\ref{sec:application}).

\subsection{Design Choice: AI Model for Gold, Price Range for All Metals}

All models in this project were trained and evaluated for gold only. Training
and validating separate models for silver, copper and platinum would repeat
the whole of Deliverable~2 three times, and the gold results suggest that the
signal would be weak for them as well (Section~\ref{sec:evaluation}). We
therefore use two levels of recommendation:

\begin{itemize}
  \item \textbf{Gold: AI model and price range.} The quote is compared with an
        expected five-day price range, and the model's predicted direction
        decides quotes that fall inside that range.
  \item \textbf{Silver, copper and platinum: price range only.} The quote is
        compared with the expected range alone. The application labels these
        results ``Price range only'' so the user knows that no AI model was
        used.
\end{itemize}

\subsection{Expected Five-Day Price Range}

The range answers the question: \emph{how far does this metal usually move in
five trading days?} For a metal with daily price $P_t$, the five-day change is
\[
  r_t = \frac{P_{t+5}}{P_t} - 1 .
\]
We compute $r_t$ for every trading day in the training years (2010 to 31
December 2021) and take its 10th, 50th and 90th percentiles,
$q_{10}$, $q_{50}$ and $q_{90}$. With the latest market price $P_{\text{last}}$,
the expected range is
\[
  \text{low} = P_{\text{last}}(1+q_{10}), \qquad
  \text{mid} = P_{\text{last}}(1+q_{50}), \qquad
  \text{high} = P_{\text{last}}(1+q_{90}).
\]
By construction, eight out of ten five-day changes in the training years fell between $q_{10}$ and $q_{90}$. This is
an empirical prediction interval built from historical changes rather than
from a fitted forecasting model~\cite{a_hyndman2021}. Only the training years
are used, so the range never contains information from the validation or test
periods, in the same way as the scaler in Deliverable~1. The application calls
the middle value the \emph{fair price} and the upper value the
\emph{walk-away price}.

Table~\ref{tab:price-ranges} shows the percentiles for the four metals. Each
is based on 3,013 five-day changes. Silver moves the most and gold the least,
so the same percentage markup can be acceptable for silver but too high for
gold.

\begin{table}[H]
  \centering
  \begin{tabular}{lcccc}
    \toprule
    Metal (unit) & $q_{10}$ (low) & $q_{50}$ (mid) & $q_{90}$ (high) & Width (points) \\
    \midrule
    Gold (USD/oz)     & $-2.5\%$ & $+0.2\%$ & $+2.7\%$ & 5.2 \\
    Silver (USD/oz)   & $-4.4\%$ & $+0.1\%$ & $+4.8\%$ & 9.2 \\
    Copper (USD/lb)   & $-3.5\%$ & $+0.1\%$ & $+3.7\%$ & 7.2 \\
    Platinum (USD/oz) & $-3.8\%$ & $-0.1\%$ & $+3.9\%$ & 7.7 \\
    \bottomrule
  \end{tabular}
  \caption{Percentiles of past five-day price changes, 2010--2021 (3,013 changes per metal).}
  \label{tab:price-ranges}
\end{table}

\subsection{Decision Rules}

Table~\ref{tab:decision-rules} lists the rules. A quote above the expected
high is always \textbf{Negotiate}, because the market rarely rises that far
in five days. When a model is used, a quote below the expected low is always
\textbf{Buy}, and inside the range the model decides: if it gives a
probability of at least 0.5 that gold will rise, waiting is likely to cost
more, so the action is \textbf{Buy}; otherwise waiting is likely to be
cheaper, so the action is \textbf{Hold}. Without a model (silver, copper and
platinum, or gold when the model files are missing), the fair price splits
the range instead.

\begin{table}[H]
  \centering
  \small
  \begin{tabular}{p{2.6cm}p{4.6cm}p{4.6cm}p{1.9cm}}
    \toprule
    Action & AI model + price range (gold) & Price range only (silver, copper, platinum) & Meaning \\
    \midrule
    \textbf{Buy} & quote $<$ low, or low $\le$ quote $\le$ high and $p_{\text{up}} \ge 0.5$ & quote $\le$ mid & buy now \\
    \textbf{Hold} & low $\le$ quote $\le$ high and $p_{\text{up}} < 0.5$ & mid $<$ quote $\le$ high & wait \\
    \textbf{Negotiate} & quote $>$ high & quote $>$ high & ask for less \\
    \bottomrule
  \end{tabular}
  \caption{Decision rules. $p_{\text{up}}$ is the selected model's probability that gold is higher in five trading days.}
  \label{tab:decision-rules}
\end{table}

\textbf{Hold} therefore has two meanings, and the application explains which
one applies. With a model it means ``the model expects the price to fall, so
waiting may be cheaper''. With the price range only it means ``the quote is
inside the usual range but above the fair price, so wait or ask for a small
discount''. The default gold model is logistic regression, because it gave
the most reliable ``up'' calls on the test years (precision 0.70,
Table~\ref{tab:test-results}). The user can switch to the random forest,
gradient boosting or the LSTM. Because the random forest and gradient
boosting predicted ``down'' on every test day, they will usually lead to
\textbf{Hold} inside the range.

\subsection{Worked Example}

With the built-in data, the last gold price is 4,153.60~USD/oz on
21 October 2026. The expected range is 4,048.39 (low), 4,162.43 (fair price)
and 4,264.93 (walk-away price). Table~\ref{tab:worked-example} shows three
quotes for 10~oz. The cost compared with the market is
$(\text{quote}-P_{\text{last}})\times\text{quantity}$.

\begin{table}[H]
  \centering
  \begin{tabular}{rrll}
    \toprule
    Quote (USD/oz) & Cost vs market (10 oz) & Action & Reason \\
    \midrule
    4,000.00 & $-1,536$ USD & Buy & below the expected low \\
    4,200.00 & $+464$ USD & Buy or Hold & inside the range: model decides \\
    4,300.00 & $+1,464$ USD & Negotiate & above the walk-away price \\
    \bottomrule
  \end{tabular}
  \caption{Gold quotes checked against the expected range (built-in data, 21 October 2026).}
  \label{tab:worked-example}
\end{table}

For the 4,200 quote the result depends on the model: it is \textbf{Buy} if
the selected model expects a rise and \textbf{Hold} if it expects a fall.
Without a model the same quote would be \textbf{Hold}, because it is above the
fair price of 4,162.43.

\subsection{Fresh Data}

The built-in data ends on 21 October 2026, so the market price it uses
becomes older every day. The user can check a quote against newer prices in
two ways: by uploading a CSV file of recent prices, or by choosing
\emph{Latest prices}, which downloads about one year of prices when the
quote is checked (Section~\ref{sec:application}). In both cases the new
prices pass through the same cleaning and feature steps as the training data.
The models themselves are not retrained, and the range percentiles still come
from 2010--2021, so the comparison stays consistent with the evaluated
models.

\subsection{Limitations}

\begin{itemize}
  \item \textbf{Exchange prices, not supplier prices.} The market price is the
        COMEX futures price. Real supplier quotes include premiums, taxes and
        delivery costs, so a fair supplier quote may sit slightly above the
        range.
  \item \textbf{Fixed width.} The range uses the same percentiles in calm and
        volatile periods. It does not widen when the market becomes more
        volatile.
  \item \textbf{Weak model signal.} The best test AUC is 0.60, so the model
        only tips quotes inside the range; it should be read as a guide, not
        a guarantee.
  \item \textbf{Buyers only.} The rules are written for a business that buys
        metal. Advice for sellers is left for future work.
\end{itemize}
